\section{Method}
\label{sec:method}

The BioShield pipeline has four interacting components: a real-text source
(PubMed), two generators, a binary detector, and a rewriting agent. They are
composed into a generic adversarial loop that any subset of components can
inhabit, which is what enables our four-condition ablation.

\subsection{Components}
\label{ssec:components}

\textbf{Real-text source.} We stream biomedical abstracts from
\texttt{ncbi/pubmed} via the Hugging Face datasets API
\cite{lhoest2021hfdatasets}. Sampling is uniform across years; abstracts
shorter than 50 tokens are discarded.

\textbf{Generator (Bio-family).} BioMistral-7B
\cite{labrak2024biomistral} is a Mistral-7B \cite{jiang2023mistral} model
continually pretrained on PubMed Central. We use it zero-shot, prompted
with biomedical topic stubs and a structured-abstract instruction. Sampling
uses temperature \(0.8\), top-\(p\) \(0.95\), and a maximum of \(384\)
new tokens.

\textbf{Generator (out-of-family baseline).} SeqGAN \cite{yu2017seqgan}
trains a small LSTM generator with policy-gradient REINFORCE against a
CNN-based discriminator. We pretrain SeqGAN on the same PubMed corpus that
BioMistral saw and use it as the Condition~B fake source.

\textbf{Detector.} We fine-tune
\texttt{microsoft/BiomedNLP-BiomedBERT-large-uncased-abstract}
\cite{gu2021biomedbert} (\(\sim\)340M parameters) as a binary
real-vs-fake classifier with a single linear head over the
\texttt{[CLS]} representation. Training uses bf16 mixed precision on
Apple~MPS, batch size \(8\), learning rate \(2 \times 10^{-5}\) with
linear decay and \(10\%\) warm-up, weight decay \(0.01\), and
\(2\) epochs per round.

\textbf{Agent (rewriter).} Qwen2.5-7B \cite{qwen2024technical} is used
zero-shot as a rewriting agent. Given a fake abstract that the current
detector found hard (probability of being real $\geq 0.5$), the agent
produces a rewrite that ``reads more like a real PubMed abstract — keep
the topic and the main claim, change the prose so it sounds like a
biomedical reviewer wrote it.'' We do not fine-tune Qwen.

\subsection{Adversarial Loop}
\label{ssec:loop}

Algorithm~\ref{alg:loop} formalises the general loop. The
\textsc{Generate} step is component~(a): SeqGAN for Condition~B,
BioMistral for Conditions~C and~D. The \textsc{HardSubset} and
\textsc{Rewrite} steps are component~(b): present in C and D, absent in
B. The \textsc{Retrain} step is component~(c): present in B and D,
absent in A and C.

\begin{algorithm}[t]
\caption{BioShield adversarial loop}
\label{alg:loop}
\begin{algorithmic}[1]
\REQUIRE Real corpus $\mathcal{R}$, generator $G$, detector $D_0$,
agent $A$ (optional), rounds $T$, fake-pool size $K$, hard-fraction $h$
\STATE $D \leftarrow D_0$  \COMMENT{baseline trained on
$(\mathcal{R}, G(\cdot))$ once}
\FOR{$r = 1, \dots, T$}
  \STATE $\textsc{set\_seed}(\text{base} + r)$
  \STATE $F_r \leftarrow \textsc{Generate}(G, K)$
  \IF{agent enabled}
    \STATE $H_r \leftarrow \textsc{HardSubset}(F_r, D, h)$
    \STATE $F_r' \leftarrow \textsc{Rewrite}(A, H_r)$
    \STATE $F_r \leftarrow (F_r \setminus H_r) \cup F_r'$
  \ENDIF
  \IF{retrain enabled}
    \STATE $D \leftarrow \textsc{Retrain}(D, \mathcal{R}, F_r)$
    \STATE $D \leftarrow \arg\max_{\text{epoch}} \mathrm{AUC}_{\text{val}}$
    \COMMENT{best-val-AUC selection}
  \ENDIF
  \STATE record $(\mathrm{AUC}, F_1, \mathrm{evasion})_r$ on test set
\ENDFOR
\RETURN $D$
\end{algorithmic}
\end{algorithm}

\subsection{Hard-Subset Selection}
We define a fake $x$ as ``hard'' for the current detector $D$ if
$D(x) \geq 0.5$, i.e., the detector misclassifies it as real. With
$h = 0.25$ and $K = 250$, the hard subset is the top-\(62\) fakes ranked by
$D$'s real-probability score. Only this subset reaches the agent — the
other \(188\) fakes are passed through unmodified. This keeps the
agent budget small (\(\sim\!8\%\) of total wall-clock,
Section~\ref{sec:setup}) while ensuring it spends compute where the
detector is actually weak.

\subsection{Checkpoint Selection}
\label{ssec:checkpoint}
After each retrain round, we keep the epoch checkpoint with the highest
validation AUC. This is the silent default in most public implementations
of adversarial-text training. We will show in Section~\ref{sec:discussion}
that this choice is what produces the C\,$=$\,D pathology: when the
baseline detector already hits the validation-AUC ceiling
(\(0.9990\) on this corpus), no retrained checkpoint can replace it, so
Condition~D's loop runs but never updates the deployed weights.

\subsection{Reproducibility Controls}
\label{ssec:repro}
We fix the random seed (\texttt{base=42}) and call
\texttt{set\_seed(42 + r)} at the top of every round so that round indices
produce different fakes and rewrites without losing reproducibility. Mixed
precision is forced to bf16 (Apple~MPS does not implement deterministic
fp16 kernels). The detector architecture, hyperparameters, validation set,
test set, agent prompt, and SeqGAN configuration are held constant across
all four conditions; the only experimental variable is which steps of
Algorithm~\ref{alg:loop} are enabled.
